\documentclass[a4paper,10pt]{article}
\usepackage[a4paper,margin=0.75in]{geometry}
\usepackage[hidelinks]{hyperref}
\usepackage{enumitem}
\usepackage{titlesec}
\usepackage{parskip}

\titleformat{\section}{\large\bfseries}{}{0em}{}[\titlerule]
\setlist[itemize]{leftmargin=*, itemsep=0pt, topsep=0pt}

\begin{document}

\begin{center}
    {\LARGE \textbf{Salaar Mir}}\\[4pt]
    London, United Kingdom \\[2pt]
    \href{mailto:salaar.mir@gmail.com}{salaar.mir@gmail.com} \quad -- \quad
    +44 7387 621 538 \quad -- \quad
    \href{https://www.linkedin.com/in/salaarmir}{linkedin.com/in/salaarmir} \quad -- \quad
    \href{https://github.com/salaarmir}{github.com/salaarmir}
\end{center}

\section*{Profile}
AI engineer with an MSc in Computational Neuroscience, Cognition \& AI and hands-on experience building Python backends, APIs, and AI-powered workflows. Strong foundation in software engineering, machine learning, and data-driven systems, with practical experience developing LLM-enabled tools for retrieval, summarisation, and decision support over real-world data. Particularly interested in agentic workflows, automation, and product-focused engineering that turns messy business processes into reliable software.

\section*{Industry Experience}
\textbf{Junior Machine Learning Engineer (Part-Time)} \hfill \textit{May 2025 -- Present}\\
\textit{Connections Tech, Remote / London}
\begin{itemize}
    \item Build Python backend services and automation pipelines for internal AI and analytics tools.
    \item Develop LLM workflows for retrieval, summarisation, and question answering over real datasets.
    \item Work with prompt design, embeddings, vector search, and structured outputs to improve quality.
    \item Design REST APIs with FastAPI / Flask and support Docker and Azure-based deployments.
    \item Collaborate with stakeholders to refine requirements and ship practical AI-enabled solutions.
\end{itemize}

\section*{Selected Projects}
\textbf{Scholar Digest -- Research Summarisation Service}
\begin{itemize}
    \item Built a Python service that collects research papers and stores structured metadata.
    \item Used LLM workflows to generate summaries and expose results through an HTTP API.
\end{itemize}

\textbf{Bayesian RL for Cognitive Flexibility (MSc Thesis)}
\begin{itemize}
    \item Implemented probabilistic reinforcement learning agents in Python.
    \item Designed experiments and analysed model behaviour using reproducible pipelines.
\end{itemize}

\textbf{Developing Maintainable Software -- Java Breakout Coursework}
\begin{itemize}
    \item Refactored and extended a Java application using modular design, Maven, JavaFX, and Git.
    \item Added JUnit testing and improved maintainability through clearer structure and documentation.
\end{itemize}

\section*{Education}
\textbf{University of Nottingham} \hfill \textit{Sept 2023 -- Dec 2024}\\
MSc Computational Neuroscience, Cognition \& AI \hfill \textit{Merit (68\%)}\\
\textit{Relevant topics:} Machine Learning in Science, Neural Computation, Practical Biomedical Modelling

\vspace{2pt}
\textbf{University of Nottingham} \hfill \textit{Sept 2020 -- July 2023}\\
BSc Computer Science with Artificial Intelligence \hfill \textit{2:1 (67\%)}\\
\textit{Relevant topics:} Machine Learning, Software Engineering, Data Structures \& Algorithms, Databases

\section*{Technical Skills}
\begin{itemize}
    \item \textbf{Languages:} Python, SQL, JavaScript, Java, C++, Bash
    \item \textbf{AI:} LLM workflows, prompt design, structured outputs, embeddings, vector search, RAG basics
    \item \textbf{Backend:} FastAPI, Flask, REST APIs, Pydantic, Git, Linux, Docker
    \item \textbf{Data:} pandas, NumPy, scikit-learn, TensorFlow, PyTorch, model evaluation
    \item \textbf{Workflow:} Azure, basic CI/CD, technical documentation, AI-assisted development with Cursor/Copilot/Claude
\end{itemize}

\end{document}
